\documentclass[10pt,a4paper]{article}

\usepackage[margin=1.8cm]{geometry}
\usepackage[T1]{fontenc}
\usepackage[utf8]{inputenc}
\usepackage{charter}
\usepackage[english]{babel}
\usepackage{enumitem}
\usepackage{hyperref}
\usepackage{xcolor}

\hypersetup{
    colorlinks=true,
    urlcolor=blue,
    linkcolor=black
}

\setlength{\parindent}{0pt}
\setlength{\parskip}{0.45em}
\setlist[itemize]{leftmargin=1.4em,topsep=2pt,itemsep=2pt}

\newcommand{\relevance}[1]{\textbf{Relevance to the Quantum Application Developer role.} #1}

\begin{document}

\begin{center}
    {\Large \textbf{Scientific Portfolio}}\\[3pt]
    {\large Antonio Michele Miti}\\[3pt]
    \textit{Selected research outputs relevant to the Quantum Application Developer position at Classiq}
\end{center}

\vspace{0.5em}

\section*{Profile}

My research portfolio combines \textbf{quantum algorithms, mathematical physics, differential geometry, and higher algebra}. 
Across these areas, a common theme is the formulation of complex physical or mathematical problems in terms of suitable
high-level structures, followed by the analysis of constraints, symmetries, reductions, and alternative representations.
For the Quantum Application Developer role, I regard this combination as particularly relevant to the development of
quantum applications: it brings together direct research experience on quantum circuits with a broader ability to
formalize technical problems, identify effective degrees of freedom, and communicate advanced concepts across disciplines.

\section*{Selected Research Outputs}

\subsection*{1. Quantum algorithms and resource-efficient arithmetic}

\textbf{A. Luongo, A. M. Miti, V. Narasimhachar, A. Sireesh,}\\
\textit{Measurement-based uncomputation of quantum circuits for modular arithmetic},\\
Proceedings of the 62nd ACM/IEEE Design Automation Conference (DAC '25), 2025.\\
DOI: \href{https://doi.org/10.1109/DAC63849.2025.11132981}{10.1109/DAC63849.2025.11132981}.\\
arXiv: \href{https://arxiv.org/abs/2407.20167}{2407.20167}.

This work investigates measurement-based uncomputation techniques for quantum circuits implementing modular arithmetic,
with particular attention to the computational resources required by arithmetic primitives and to the cost of reversing
intermediate computations.

\relevance{This is the most direct connection between my research and Classiq's activity. It demonstrates research
experience at the quantum algorithm/circuit level, including uncomputation, modular arithmetic, circuit-resource
trade-offs, and the analysis of implementation constraints. These are closely related to the problems encountered when
designing, synthesizing, and optimizing quantum applications at a high level.}

\subsection*{2. Reduction of constrained mathematical systems}

\textbf{A. M. Miti, L. Ryvkin,}\\
\textit{Multisymplectic observable reduction using constraint triples},\\
Differential Geometry and its Applications \textbf{100} (2025), 102272.\\
DOI: \href{https://doi.org/10.1016/j.difgeo.2025.102272}{10.1016/j.difgeo.2025.102272}.\\
arXiv: \href{https://arxiv.org/abs/2506.00234}{2506.00234}.

The paper develops a systematic framework for reducing observables in multisymplectic systems subject to constraints,
combining geometric structure with an explicit treatment of the data required for reduction.

\relevance{Although the setting is mathematical physics rather than quantum computing, the underlying methodology is
highly transferable: identify constraints and symmetries, isolate the effective structure, and construct a reduced
description that preserves the information relevant to the problem. This type of abstraction is useful when translating
scientific problems into compact and efficient quantum formulations.}

\subsection*{3. Structure-preserving simplification of higher algebraic models}

\textbf{C. Blacker, A. M. Miti, L. Ryvkin,}\\
\textit{Reduction of $L_\infty$-algebras of observables on multisymplectic manifolds},\\
SIGMA \textbf{20} (2024), 061.\\
DOI: \href{https://doi.org/10.3842/SIGMA.2024.061}{10.3842/SIGMA.2024.061}.\\
arXiv: \href{https://arxiv.org/abs/2206.03137}{2206.03137}.

This work studies how a complex algebraic model of physical observables behaves under symmetry reduction, with emphasis
on retaining the algebraic relations that encode the original system.

\relevance{The work demonstrates experience in decomposing and simplifying complex mathematical structures while
preserving their essential relations. For an application-oriented quantum role, this supports the ability to reason about
high-level representations, dependencies, constraints, and transformations before moving to a lower-level implementation.}

\subsection*{4. High-level modelling of observables and alternative representations}

\textbf{A. M. Miti, M. Zambon,}\\
\textit{Observables on multisymplectic manifolds and higher Courant algebroids},\\
Communications in Contemporary Mathematics \textbf{27} (2025), no.~8, 2550006.\\
DOI: \href{https://doi.org/10.1142/S0219199725500063}{10.1142/S0219199725500063}.\\
arXiv: \href{https://arxiv.org/abs/2209.05836}{2209.05836}.

The paper compares and connects different algebraic descriptions of observables arising in higher geometric mechanics,
clarifying how information can be encoded in distinct but related mathematical structures.

\relevance{This work reflects a recurring aspect of application development: choosing an appropriate representation of
a problem and understanding how different representations encode the same underlying information. It also demonstrates
the ability to work with sophisticated abstractions and explain the relation between them rigorously.}

\subsection*{5. Mappings between complex algebraic structures}

\textbf{D. Fiorenza, A. M. Miti,}\\
\textit{$L_\infty$-morphisms between twisted Courant $r$-Lie algebras and untwisted Courant $(r+1)$-Lie algebroids},\\
preprint, 2026.\\
arXiv: \href{https://arxiv.org/abs/2602.14702}{2602.14702}.

This work constructs and studies explicit higher-algebraic mappings between two different formulations of related
geometric structures.

\relevance{Its relevance lies primarily in the methodology: translating between abstract representations while keeping
track of the structural information that must be preserved. This is complementary to my quantum-computing work and
illustrates the level of mathematical abstraction and formal reasoning I can bring to application development.}

\section*{Additional Peer-Reviewed Research}

My broader publication record establishes a sustained background in mathematical physics, symmetry, quantization, and
geometric modelling. Selected additional papers include:

\begin{itemize}
    \item \textbf{A. M. Miti, L. Ryvkin,}
    \textit{Multisymplectic actions of compact Lie groups on spheres},
    Journal of Symplectic Geometry \textbf{18} (2020), no.~6, 1751--1785.
    DOI: \href{https://doi.org/10.4310/JSG.2020.v18.n6.a6}{10.4310/JSG.2020.v18.n6.a6}.

    \item \textbf{A. M. Miti, M. Spera,}
    \textit{Derivation of the HOMFLYPT knot polynomial via helicity and geometric quantization},
    Bollettino dell'Unione Matematica Italiana \textbf{14} (2021), 269--284.
    DOI: \href{https://doi.org/10.1007/s40574-020-00254-5}{10.1007/s40574-020-00254-5}.

    \item \textbf{A. M. Miti, M. Spera,}
    \textit{A Hydrodynamical Homotopy Co-momentum Map and a Multisymplectic Interpretation of Higher-order Linking Numbers},
    Journal of the Australian Mathematical Society (2021).
    DOI: \href{https://doi.org/10.1017/S1446788720000518}{10.1017/S1446788720000518}.
\end{itemize}

These works are less directly connected to quantum software, but they document the breadth of my experience in
\textbf{mathematical modelling, symmetries, conserved quantities, quantization, and collaborative research}.

\section*{Complementary Technical Perspective}

The scientific portfolio is complemented by experience with \textbf{C++, Python, CUDA, Linux, parallel computing, and
scientific software}. I have also worked on GPU-oriented implementations and performance analysis of mathematical
algorithms. Together with my quantum-circuit research, this gives me a profile that connects theoretical modelling with
the computational considerations involved in implementing advanced scientific applications.

My academic activity has additionally involved teaching, research seminars, international collaborations, and the
presentation of technically demanding material to audiences with different backgrounds. I therefore see the
Quantum Application Developer role as a natural setting in which to combine \textbf{quantum algorithms, scientific
computing, mathematical abstraction, and technical communication}.

\vfill

\noindent
\href{https://www.antoniomiti.it/}{www.antoniomiti.it}
\hfill
\href{https://orcid.org/0000-0002-8829-1943}{ORCID: 0000-0002-8829-1943}

\end{document}
